
设 $(M,\omega)$ 是闭\BookRecordIndex{T0358}辛流形，$J\in\mathcal J_\tau(M,\omega)$ 是被 $\omega$ \BookRecordIndex{T0372}驯服的
\BookRecordIndex{T0088}殆复结构。固定一个球面同调类 $A\in H_2(M)$ 和一个非负整数 $k$。当 $A\ne0$ 时，记
\[
  \mathcal M^*_{0,k}(A;J)
  :=\frac{\mathcal M^*(A;J)\times\bigl((S^2)^k\setminus\Delta_k\bigr)}{G}
\]
为元组 $(u,z_1,\ldots,z_k)$ 的等价类所成的\BookRecordIndex{T0246}模空间，其中
$u:S^2\to M$ 是表示类 $A$ 的单纯 \BookRecordIndex{T0003}$J$-全纯曲线，$z_i$ 是 $S^2$ 上两两不同的点。
记号 $\mathcal M^*(A;J)$ 见第~\ref{sec:3-1} 节；
$\Delta_k\subset(S^2)^k$ 是\BookRecordIndex{T0082}大对角线，即由 $S^2$ 中至少两个点重合的有序 $k$-元组组成的
集合。等价关系由\BookRecordIndex{T0413}重参数化群 $G=\operatorname{PSL}(2,\mathbb C)$ 的对角作用给出。当
$k\ge3$ 时，我们约定常值映射也是单纯的，从而把这一记号推广到类 $A=0$；当 $k<3$
时，则定义 $\mathcal M^*_{0,k}(0;J)$ 为空集。略微滥用语言，即使 $A=0$，我们仍称
$\mathcal M^*_{0,k}(A;J)$ 为\emph{类 $A$ 中带 $k$ 个\BookRecordIndex{T0066}标记点的单纯 $J$-全纯球面的\BookRecordIndex{T0246}模空间}。
用第~\ref{sec:5-1} 节的记号，它是表示类 $A$、带 $k$ 个\BookRecordIndex{T0066}标记点的亏格零\BookRecordIndex{T0340}稳定映射\BookRecordIndex{T0246}模空间
$\overline{\mathcal M}_{0,k}(A;J)$ 的\BookRecordIndex{T0106}顶层 $\mathcal M_{0,k}(A;J)$ 的一个开子集。
由第 3 章的结果，对通有的\BookRecordIndex{T0088}殆复结构 $J$，\BookRecordIndex{T0246}模空间 $\mathcal M^*_{0,k}(A;J)$ 是一个
（非紧）光滑流形，其维数为
\[
  \dim\mathcal M^*_{0,k}(A;J)
  =\mu(A,k):=2n+2c_1(A)+2k-6.
\]
本章的目标是给出条件，使\BookRecordIndex{T0283}求值映射
\[
  \operatorname{ev}:\mathcal M^*_{0,k}(A;J)\longrightarrow M^k,
  \qquad
  \operatorname{ev}([u,z_1,\ldots,z_k])
  :=\bigl(u(z_1),\ldots,u(z_k)\bigr)
\]
在 $M^k$ 中定义同调类
$\operatorname{ev}_*([\overline{\mathcal M}_{0,k}(A;J)])$。第 7 章将用这个类定义
\BookRecordIndex{T0143}格罗莫夫--威滕不变量。

直观地说，若映射 $f:V\to X$ 的定义域是一个非紧 $d$-维流形，而它的极限集
$f(\overline V\setminus V)$ 的维数至多为 $d-2$，那么 $f$ 在 $X$ 中表示一个 $d$ 维
同调类。这里，$\overline V$ 应理解为 $V$ 的某个适当紧化，$f$ 可以连续延拓到其上；维数条件
则是说，“边界” $\partial V:=\overline V\setminus V$ 在 $f$ 下的像包含在某个光滑映射
的像中，而该映射定义域的每个连通分支维数都不超过 $d-2$。我们称这样的映射
$f:V\to X$ 为一个 $d$ 维\emph{\BookRecordIndex{T0337}伪循环}；精确定义见第~\ref{sec:6-5} 节。为了验证\BookRecordIndex{T0283}求值映射
$\operatorname{ev}:\mathcal M^*_{0,k}(A;J)\to M^k$ 的\BookRecordIndex{T0337}伪循环条件，必须考察边界
\[
  \overline{\mathcal M}_{0,k}(A;J)\setminus\mathcal M^*_{0,k}(A;J)
\]
在延拓后的\BookRecordIndex{T0283}求值映射下的像。这个边界包含两类元素：
$\mathcal M_{0,k}(A;J)$ 中非单纯的元素，以及属于满足 $e(T)>0$ 的\BookRecordIndex{T0072}层
$\mathcal M_{0,T}(A;J)$ 的\BookRecordIndex{T0340}稳定映射；这里 $e(T)$ 是树 $T$ 的边数。我们将看到，每个\BookRecordIndex{T0072}层
$\mathcal M_{0,T}(A;J)$ 的虚维数都是 $\mu(A,k)-2e(T)$。因此，若能忽略非单纯元素，
并通过选取通有的 $J$ 使每个\BookRecordIndex{T0072}层正则化，那么所需的维数条件便会成立。

当然，一般而言这是做不到的。正如第 3 章一样，关于通有\BookRecordIndex{T0088}殆复结构 $J$ 的\BookRecordIndex{T0166}横截性结果，只能在
限制到\BookRecordIndex{T0184}简单稳定映射的子空间 $\mathcal M^*_{0,T}(\{B_\alpha\};J)$ 后建立。
命题~\ref{prop:6-1-2} 断言，只要考虑所有满足
\[
  A=\sum_{\alpha\in T}m_\alpha B_\alpha,
  \qquad m_\alpha>0,
\]
的类族 $\{B_\alpha\}_{\alpha\in T}$，就确实可以只用\BookRecordIndex{T0183}简单层
$\mathcal M^*_{0,T}(\{B_\alpha\};J)$ 来紧化 $\mathcal M^*_{0,k}(A;J)$。现在的问题是，
一般并不能保证\BookRecordIndex{T0072}层 $\mathcal M^*_{0,T}(\{B_\alpha\};J)$ 的维数
$\mu(B,k)-2e(T)$ 足够小，即不超过 $\mu(A,k)-2$。若要它对所有类 $A$ 都成立，就必须假设
$(M,\omega)$ \BookRecordIndex{T0058}半正；这意味着对通有的 $J$，不存在陈数为负的全纯球面。它是代数几何中
NEF 条件的辛版本。当 $e(T)=0$ 时，还必须谨慎处理陈数为零的全纯球面。

最终结果是：当 $(M,\omega)$ \BookRecordIndex{T0058}半正时，\BookRecordIndex{T0283}求值映射
\[
  \operatorname{ev}:\mathcal M^*_{0,k}(A;J)\longrightarrow M^k
\]
是一个\BookRecordIndex{T0337}伪循环，称为\emph{\BookRecordIndex{T0145}格罗莫夫--威滕伪循环}。我们将在第~\ref{sec:6-6} 节证明这一点。
证明的主要工具是定理~\ref{thm:6-2-6}；它断言，对通有的 $J$，\BookRecordIndex{T0246}模空间
$\mathcal M^*_{0,T}(\{B_\alpha\};J)$ 是有限维流形，且
\[
  \dim\mathcal M^*_{0,T}(\{B_\alpha\};J)
  =\mu(B,k)-2e(T),
  \qquad B:=\sum_{\alpha\in T}B_\alpha.
\]
若 $(M,\omega)$ \BookRecordIndex{T0058}半正，则对每个满足
$\sum_{\alpha\in T}m_\alpha B_\alpha=A$、$m_\alpha>0$ 的同调类族
$\{B_\alpha\}_\alpha$，以及每棵满足 $e(T)>0$ 的树 $T$，这个数都不超过
$\mu(A,k)-2$。当 $e(T)=0$ 时，只有在 $c_1(A)>0$ 时才会得到低维\BookRecordIndex{T0246}模空间。若
$A=mB\in H_2(M)$（$m>1$）是陈数为零的类的倍数，就需要更广的一类扰动来控制
$\mathcal M^*_{0,k}(A;J)$ 的边界。我们将在第~\ref{sec:6-7} 节说明，此时可以定义\BookRecordIndex{T0337}伪循环
\[
  \operatorname{ev}:\mathcal M^*_{0,k}(A;\{J_z\})\longrightarrow M^k,
\]
不过定义方程中的\BookRecordIndex{T0088}殆复结构 $J_z$ 现在依赖于基点 $z\in S^2$。这一办法还使我们能够定义
某些\BookRecordIndex{T0066}标记点固定的\BookRecordIndex{T0337}伪循环。

本章最后简要讨论\BookRecordIndex{T0283}求值映射与\BookRecordIndex{T0377}遗忘映射 $\pi$ 的乘积：
\[
  \operatorname{ev}\times\pi:
  \overline{\mathcal M}_{0,k}(A;J)\longrightarrow M^k\times\overline{\mathcal M}_{0,k}.
\]
上面提到的所有\BookRecordIndex{T0337}伪循环，无论\BookRecordIndex{T0066}标记点固定还是可变，都可以通过限制这一乘积映射，再投影到第一个
因子而得到。然而，对任意\BookRecordIndex{T0058}半正流形证明 $\operatorname{ev}\times\pi$ 本身是\BookRecordIndex{T0337}伪循环，需要
更广的一类扰动；因此，我们只处理一个很特殊的情形，见习题 6.7.13。

要为一般紧\BookRecordIndex{T0358}辛流形构造\BookRecordIndex{T0143}格罗莫夫--威滕不变量，就必须正则化整个紧化
$\overline{\mathcal M}_{0,k}(A;J)$，包括其中的\BookRecordIndex{T0113}多重覆盖曲线。这会导向格罗莫夫--威滕理论中
\BookRecordIndex{T0369}虚基本循环的构造；自二十世纪九十年代中期以来，这一直是大量研究的主题（见
Fukaya--Ono\,\srcite{127}、Li--Tian\,\srcite{239}、Ruan\,\srcite{343}、
Liu--Tian\,\srcite{249}、Siebert\,\srcite{376}、McDuff\,\srcite{267}）。
在近期一系列论文中，Hofer--Wysocki--Zehnder 用\BookRecordIndex{T0076}尺度微积分和 \BookRecordIndex{T0044}polyfold 重新表述了整套理论，
见 \srcite{184,185,186,187,188,189}。基于本书所讨论分析方法的较传统处理，可见
McDuff--Wehrheim\,\srcite{281}。本书不讨论\BookRecordIndex{T0369}虚基本循环，而只关注\BookRecordIndex{T0058}半正情形；在那里，可以
直接通过计数 \BookRecordIndex{T0003}$J$-全纯曲线来定义\BookRecordIndex{T0143}格罗莫夫--威滕不变量。

与前几章一样，我们的大多数论证都限于亏格零情形。不过，当 $u$ 的定义域
$(\Sigma,j_\Sigma)$ 是复结构固定的闭黎曼曲面时，其中有些论证经小幅修改仍然成立。

\section{\BookRecordIndex{T0184}简单稳定映射}\label{sec:6-1}
\zhindex{jiandanwendingyingshe}{简单稳定映射}

本节引入\BookRecordIndex{T0184}简单稳定映射的概念，并为每个\BookRecordIndex{T0340}稳定映射构造一个与它有相同像的\BookRecordIndex{T0184}简单稳定映射。
正是这个关键性质，使我们可以只用\BookRecordIndex{T0183}简单层对\BookRecordIndex{T0246}模空间 $\mathcal M^*_{0,k}(A;J)$ 作有用的紧化。

\begin{definition}\label{def:6-1-1}
设\BookRecordIndex{T0340}稳定映射
\[
  (\mathbf u,\mathbf z)\in\mathcal{SC}_{0,k}(M,A;J)
\]
以一棵 $k$-\BookRecordIndex{T0084}带标号树 $T=(T,E,\Lambda)$ 为模型。若每个非常值映射 $u_\alpha$ 都是单纯的，
并且对任意两个不同顶点 $\alpha\ne\beta$，只要 $u_\alpha,u_\beta$ 都非常值，就有
\[
  u_\alpha(S^2)\ne u_\beta(S^2),
\]
则称该\BookRecordIndex{T0340}稳定映射是\emph{简单的}。
\end{definition}

设 $T=(T,E,\Lambda)$ 是一棵 $k$-\BookRecordIndex{T0084}带标号树，并设函数
\[
  T\longrightarrow H_2(M;\mathbb Z),
  \qquad \alpha\longmapsto A_\alpha
\]
满足
\begin{equation}\label{eq:6-1-1}
  A_\alpha=0
  \quad\Longrightarrow\quad
  \#\bigl(\Lambda_\alpha\cup\{\beta\in T\mid\alpha E\beta\}\bigr)\ge3.
\end{equation}
用 $\widetilde{\mathcal M}^*_{0,T}(\{A_\alpha\};J)$ 表示所有以 $T$ 为模型、且对每个
$\alpha\in T$ 都满足 $[u_\alpha]=A_\alpha$ 的\BookRecordIndex{T0184}简单稳定映射
$(\mathbf u,\mathbf z)$ 所成的集合。由定义容易看出，\BookRecordIndex{T0413}重参数化群 $G_T$ 在
$\widetilde{\mathcal M}^*_{0,T}(\{A_\alpha\};J)$ 上的作用自由且恰当；记其商为
\[
  \mathcal M^*_{0,T}(\{A_\alpha\};J)
  :=\widetilde{\mathcal M}^*_{0,T}(\{A_\alpha\};J)/G_T.
\]
当 $T$ 没有边时，这个商正是本章导言中提到的带 $k$ 个\BookRecordIndex{T0066}标记点的单纯 $J$-全纯球面空间
$\mathcal M^*_{0,k}(A;J)$。由稳定性条件，当 $A=0$ 且 $k<3$ 时，这个空间为空。

任意\BookRecordIndex{T0340}稳定映射可能由于两种原因之一而不是简单的：它可能含有某些多重覆盖分支，或者含有若干
像相同的分支。本节余下部分将证明，每个\BookRecordIndex{T0340}稳定映射都覆盖一个与它具有相同像的\BookRecordIndex{T0184}简单稳定映射。
直观上这很清楚：只需删去像相同的分支，再把每个多重覆盖替换成其底\BookRecordIndex{T0072}层单纯映射即可。然而，
还须证明所得\BookRecordIndex{T0340}稳定映射仍以一棵树为模型，并且还必须追踪\BookRecordIndex{T0066}标记点。下一个命题给出详细说明。
回忆 $\operatorname{ev}_i$ 表示在第 $i$ 个\BookRecordIndex{T0066}标记点处求值，见习题~\ref{ex:5-1-11}。

\begin{proposition}\label{prop:6-1-2}
设
$(\mathbf u,\mathbf z)\in\widetilde{\mathcal M}_{0,T}(\{A_\alpha\};J)$
\sourcefootnote{原书把命题的输入空间印作带星号的
$\widetilde{\mathcal M}^{*}_{0,T}$；但本节刚把星号空间定义为\BookRecordIndex{T0184}简单稳定映射空间，而命题前文、
结论及以下三步证明均从任意\BookRecordIndex{T0340}稳定映射出发，故去掉星号。}
是一个以 $k$-\BookRecordIndex{T0084}带标号树 $(T,E,\Lambda)$ 为模型、并对 $\alpha\in T$ 表示类
$A_\alpha\in H_2(M;\mathbb Z)$ 的\BookRecordIndex{T0340}稳定映射。则存在一棵 $k$-\BookRecordIndex{T0084}带标号树
$T'=(T',E',\Lambda')$、两个函数
\[
  T'\longrightarrow H_2(M;\mathbb Z):\alpha'\longmapsto A'_{\alpha'},
  \qquad
  T'\longrightarrow\mathbb N:\alpha'\longmapsto m_{\alpha'},
\]
以及\BookRecordIndex{T0184}简单稳定映射
$(\mathbf u',\mathbf z')\in
\widetilde{\mathcal M}^*_{0,T'}(\{A'_{\alpha'}\};J)$，使得
\[
  \bigcup_{\alpha'\in T'}u'_{\alpha'}(S^2)
  =\bigcup_{\alpha\in T}u_\alpha(S^2),
  \qquad
  A=\sum_{\alpha'\in T'}m_{\alpha'}A'_{\alpha'},
\]
并且对 $i=1,\ldots,k$，
\[
  \operatorname{ev}_i(\mathbf u',\mathbf z')
  =\operatorname{ev}_i(\mathbf u,\mathbf z).
\]
\end{proposition}

\begin{proof}
我们证明 $k=0$，即没有\BookRecordIndex{T0066}标记点的情形；一般情形留作习题。定义一个（没有\BookRecordIndex{T0066}标记点的）
\emph{\BookRecordIndex{T0180}加权预稳定映射}为元组
\[
  (T,f,\mathbf m,\mathbf u,\mathbf z)
  =\bigl(T,f,\{m_\alpha\}_{\alpha\in T},
    \{u_\alpha\}_{\alpha\in T},
    \{z_\alpha\}_{\alpha\in T\setminus\{0\}}\bigr),
\]
其中 $(T,\preceq)$ 是最小元为 $0$ 的有限全序集，
$f:T\setminus\{0\}\to T$ 是一个函数（它起着习题~\ref{ex:5-1-3} 中向量
$\mathbf j$ 的作用），$m_\alpha\in\mathbb N$ 是正整数权，
$u_\alpha:S^2=\mathbb C\cup\{\infty\}\to M$ 是 $J$-全纯球面，
$z_\alpha\in\mathbb C$ 是复数，并且对每个 $\alpha\in T\setminus\{0\}$ 都满足
\[
  f(\alpha)\prec\alpha,
  \qquad
  u_{f(\alpha)}(z_\alpha)=u_\alpha(\infty).
\]
进一步假设，只要 $u_\alpha$ 是常值映射，集合
\[
  T_\alpha:=\{\beta\in T\mid\alpha\prec\beta, f(\beta)=\alpha\}
\]
就至少含两个点。可把这些数据看成对一棵 $J$-全纯球面树的归纳描述：每个新球面
$u_\alpha$ 都在点
$u_\alpha(\infty)=u_{f(\alpha)}(z_\alpha)$ 处连接到先前的某个球面，即
$u_{f(\alpha)}$。因此，应把点 $z_\alpha$ 看作属于由 $f(\alpha)$ 标号的球面。数
$m_\alpha$ 被理解为 $u_\alpha$ 的重数，而\BookRecordIndex{T0180}加权预稳定映射
$(T,f,\mathbf m,\mathbf u,\mathbf z)$ 所表示的同调类定义为
\[
  A:=\sum_{\alpha\in T}m_\alpha A_\alpha,
  \qquad A_\alpha:=[u_\alpha].
\]
预\BookRecordIndex{T0340}稳定映射不必稳定，因为并未要求同一球面上的\BookRecordIndex{T0191}节点两两不同。

若对所有 $\alpha,\beta\in T\setminus\{0\}$ 都有
\[
  \alpha\ne\beta,\quad f(\alpha)=f(\beta)
  \quad\Longrightarrow\quad z_\alpha\ne z_\beta,
\]
则称\BookRecordIndex{T0180}加权预稳定映射 $(T,f,\mathbf m,\mathbf u,\mathbf z)$ 是\emph{稳定的}。若每当
$u_\alpha$ 非常值时它都是单纯的，而且每当 $\alpha\ne\beta$ 且
$u_\alpha,u_\beta$ 都非常值时都有
$u_\alpha(S^2)\ne u_\beta(S^2)$，则称它是\emph{简单的}。我们分三步证明：对每个\BookRecordIndex{T0180}加权
预稳定映射 $(T,f,\mathbf m,\mathbf u,\mathbf z)$，都存在一个与它有相同像、表示相同同调类
的简单加权\BookRecordIndex{T0340}稳定映射 $(T',f',\mathbf m',\mathbf u',\mathbf z')$。

\emph{步骤 1。}
给定\BookRecordIndex{T0180}加权预稳定映射 $(T,f,\mathbf m,\mathbf u,\mathbf z)$，存在一个加权\BookRecordIndex{T0340}稳定映射
$(T',f',\mathbf m',\mathbf u',\mathbf z')$ 和一个单射 $\iota:T\to T'$，使得对每个
$\alpha\in T$，
\[
  u_\alpha=u'_{\iota(\alpha)},
  \qquad m_\alpha=m'_{\iota(\alpha)},
\]
而对每个 $\alpha'\in T'\setminus\iota(T)$，$u'_{\alpha'}$ 都是常值映射。特别地，两个映射
具有相同的像，并表示相同的同调类。

设存在顶点 $\alpha_0,\alpha_1\in T\setminus\{0\}$，使得
\[
  \alpha_0\ne\alpha_1,
  \qquad f(\alpha_0)=f(\alpha_1)=:\gamma,
  \qquad z_{\alpha_0}=z_{\alpha_1}.
\]
在 $\gamma$ 的紧继位置插入一个新元素 $\gamma'$，定义新树
$T':=T\cup\{\gamma'\}$。新次序为
\[
  \alpha\prec'\beta
  \quad\Longleftrightarrow\quad
  \begin{cases}
    \alpha\prec\beta, & \alpha,\beta\in T,\\
    \alpha\preceq\gamma, & \beta=\gamma',\ \alpha\in T,\\
    \gamma\prec\beta, & \alpha=\gamma',\ \beta\in T.
  \end{cases}
\]
新的权和全纯曲线定义为
\[
  m'_\alpha:=
  \begin{cases}
    m_\alpha,&\alpha\in T,\\
    1,&\alpha=\gamma',
  \end{cases}
  \qquad
  u'_\alpha:=
  \begin{cases}
    u_\alpha,&\alpha\in T,\\
    u_\gamma(z_{\alpha_0}),&\alpha=\gamma'.
  \end{cases}
\]
函数 $f':T'\setminus\{0\}\to T'$ 与\BookRecordIndex{T0191}节点定义为
\[
  f'(\alpha):=
  \begin{cases}
    f(\alpha),&\alpha\in T\setminus\{\alpha_0,\alpha_1\},\\
    \gamma',&\alpha\in\{\alpha_0,\alpha_1\},\\
    \gamma,&\alpha=\gamma',
  \end{cases}
  \qquad
  z'_\alpha:=
  \begin{cases}
    z_\alpha,&\alpha\in T\setminus\{\alpha_0,\alpha_1\},\\
    0,&\alpha=\alpha_0,\\
    1,&\alpha=\alpha_1,\\
    z_{\alpha_0},&\alpha=\gamma'.
  \end{cases}
\]
如此归纳下去，便能构造出与原预\BookRecordIndex{T0340}稳定映射有相同像的加权\BookRecordIndex{T0340}稳定映射。这证明了步骤 1。

\emph{步骤 2。}
给定\BookRecordIndex{T0180}加权预稳定映射 $(T,f,\mathbf m,\mathbf u,\mathbf z)$，存在一个与它有相同像、表示相同
同调类的加权\BookRecordIndex{T0340}稳定映射 $(T',f',\mathbf m',\mathbf u',\mathbf z')$，使得每个非常值
$u'_{\alpha'}$ 都是单纯的。

设某个 $\gamma\in T$ 对应的 $u_\gamma$ 非常值且为多重覆盖。由命题~\ref{prop:2-5-1}，
存在单纯 $J$-全纯球面 $v:S^2\to M$ 和次数
$d:=\deg(\phi)>1$ 的有理映射 $\phi:S^2\to S^2$，使得
$u_\gamma=v\circ\phi$ 且 $\phi(\infty)=\infty$。这给出一个\BookRecordIndex{T0180}加权预稳定映射
$(T,f,\mathbf m',\mathbf u',\mathbf z')$，其中对 $\alpha\in T$，
\[
  m'_\alpha:=
  \begin{cases}
    m_\alpha,&\alpha\ne\gamma,\\
    dm_\gamma,&\alpha=\gamma,
  \end{cases}
  \qquad
  u'_\alpha:=
  \begin{cases}
    u_\alpha,&\alpha\ne\gamma,\\
    v,&\alpha=\gamma,
  \end{cases}
\]
而对 $\alpha\in T\setminus\{0\}$，
\[
  z'_\alpha:=
  \begin{cases}
    z_\alpha,&f(\alpha)\ne\gamma,\\
    \phi(z_\alpha),&f(\alpha)=\gamma.
  \end{cases}
\]
应用步骤 1 使映射稳定，再归纳继续。这证明了步骤 2。

\emph{步骤 3。}
命题在 $k=0$ 时成立。

设 $(T,f,\mathbf m,\mathbf u,\mathbf z)$ 是加权\BookRecordIndex{T0340}稳定映射，且每个非常值 $u_\alpha$ 都是
单纯的。设存在顶点 $\gamma_0,\gamma_1\in T$，满足
\[
  \gamma_0\prec\gamma_1,
  \qquad
  u_{\gamma_0}(S^2)=u_{\gamma_1}(S^2),
\]
且 $u_{\gamma_0},u_{\gamma_1}$ 都非常值。由推论~\ref{cor:2-5-4}，存在 \BookRecordIndex{T0040}Möbius 变换
$\phi:S^2\to S^2$，使
\[
  u_{\gamma_1}=u_{\gamma_0}\circ\phi.
\]
令 $T':=T\setminus\{\gamma_1\}$，赋予它由 $(T,\prec)$ 诱导的次序，并对
$\alpha\in T'$ 令 $u'_\alpha:=u_\alpha$。对 $\alpha\in T'\setminus\{0\}$ 定义
\[
  f'(\alpha):=
  \begin{cases}
    f(\alpha),&f(\alpha)\ne\gamma_1,\\
    \gamma_0,&f(\alpha)=\gamma_1,
  \end{cases}
  \qquad
  z'_\alpha:=
  \begin{cases}
    z_\alpha,&f(\alpha)\ne\gamma_1,\\
    \phi(z_\alpha),&f(\alpha)=\gamma_1,
  \end{cases}
\]
并对 $\alpha\in T'$ 定义
\[
  m'_\alpha:=
  \begin{cases}
    m_\alpha,&\alpha\ne\gamma_0,\\
    m_{\gamma_0}+m_{\gamma_1},&\alpha=\gamma_0.
  \end{cases}
\]
然后用步骤 1 使映射稳定，再归纳继续。这便证明了命题在 $k=0$ 时成立。
\end{proof}

\begin{exercise}\label{ex:6-1-3}
证明命题~\ref{prop:6-1-2} 在 $k>0$ 时成立。提示：还须在\BookRecordIndex{T0184}简单稳定映射上选取\BookRecordIndex{T0066}标记点，使其像
与原曲线的\BookRecordIndex{T0066}标记点之像相同。这可能需要引入额外的（幽灵）顶点。
\end{exercise}

\begin{exercise}\label{ex:6-1-4}
证明\BookRecordIndex{T0413}重参数化群 $G_T$ 在空间
$\widetilde{\mathcal M}^*_{0,T}(\{A_\alpha\};J)$ 上的作用自由且恰当。
\end{exercise}
